\section{Translation-invariant Hamiltonians on a finite set of sites}
\label{sec:hastings_site_symmetry}

% Hastings, arXiv:cond-mat/0508554v3, proj3new.tex, lines 862--863 and
% 1603--1642.

\cite{Hastings2006Solving} shows that for a translation-invariant system on a
lattice the low-energy commutators of the filtered terms are exponentially small
in the linear size of the system. The statements here describe translation
invariance on any finite set of sites $V$, with a group $G$ acting on $V$, such as
the translations of a periodic lattice or a group of automorphisms of a finite
graph. Each $g\in G$ moves the configurations $x\colon V\to\{1,\dots,q\}$ to
$g\cdot x=x\circ g^{-1}$, and $P_g$ is the permutation matrix of this move.

\begin{definition}[Site permutations]
    \label{def:hastings_site_permutation}
    \lean{Matrix.sitePermMatrix, Matrix.sitePermMatrix_mul,
        Matrix.sitePermMatrix_mem_unitaryGroup,
        Matrix.sitePermMatrix_mul_mul_conjTranspose}
    \leanok
    For $g\in G$, $P_g$ is the permutation matrix of $x\mapsto g\cdot x$ on the
    configurations. The map $g\mapsto P_g$ is a unitary representation of $G$, and
    $(P_gAP_g^\dagger)(x,y)=A(x\circ g,y\circ g)$ for every operator $A$.
\end{definition}

\begin{theorem}[Support of a translated operator]
    \label{thm:hastings_site_permutation_support}
    \lean{Matrix.sitePermMatrix_mul_mul_conjTranspose_mem_supportedOperators}
    \leanok
    \uses{def:hastings_site_permutation,aux:permuted-operator-support}
    If $A$ acts on the sites $S$, then $P_gAP_g^\dagger$ acts on $g\cdot S$.
\end{theorem}
\begin{proof}\leanok
    $P_gAP_g^\dagger$ is $A$ with both configurations precomposed with the
    permutation $v\mapsto g\cdot v$ of the sites, which carries the support $S$ to
    $g\cdot S$.
\end{proof}

\begin{definition}[Invariant operators]
    \label{def:hastings_site_invariant}
    \lean{Matrix.IsSiteInvariant, Matrix.isSiteInvariant_of_conj}
    \leanok
    \uses{def:hastings_site_permutation}
    An operator $H$ is invariant under the action of $G$ when $P_gH=HP_g$ for every
    $g\in G$. In particular, $H$ is invariant if $P_gHP_g^\dagger=H$ for every
    $g\in G$.
\end{definition}

\begin{theorem}[Sums of translated terms]
    \label{thm:hastings_site_invariant_sum}
    \lean{Matrix.isSiteInvariant_sum_of_covariant, Matrix.isSiteInvariant_sum_conj}
    \leanok
    \uses{def:hastings_site_invariant}
    Let $G$ also act on a finite index set $K$, and let the operators $h_k$ satisfy
    $h_{g\cdot k}=P_gh_kP_g^\dagger$. Then $\sum_kh_k$ is invariant. In particular,
    for one operator $H_0$, the sum $\sum_{g\in G}P_gH_0P_g^\dagger$ of its
    translates is invariant.
\end{theorem}
\begin{proof}\leanok
    \uses{def:hastings_site_permutation}
    Conjugating by $P_g$ sends $h_k$ to $h_{g\cdot k}$, and $k\mapsto g\cdot k$
    permutes $K$. For the translates, take $K=G$ acting on itself by left
    multiplication.
\end{proof}

\begin{theorem}[Filtering commutes with translation]
    \label{thm:hastings_filtered_term_translate}
    \lean{TNLean.Lattice.gaussianFilteredTerm_conj,
        TNLean.Lattice.gaussianFilteredTerm_sitePermMatrix_conj}
    \leanok
    \uses{def:hastings_site_invariant,def:hastings_gaussian_filtered_term}
    For a unitary $U$ and Hermitian $H$, filtering $UWU^\dagger$ with respect to
    $UHU^\dagger$ gives $U\widetilde WU^\dagger$. If $H$ is invariant, the filtered
    operator of $P_gWP_g^\dagger$ is therefore $P_g\widetilde WP_g^\dagger$.
\end{theorem}
\begin{proof}\leanok
    \uses{def:hastings_site_permutation}
    $e^{iUHU^\dagger t}=Ue^{iHt}U^\dagger$, so the integrand of the filter is
    conjugated by $U$ at every time, and conjugation by $U$ commutes with the
    integral.
\end{proof}

\begin{theorem}[Graph automorphisms preserve distances]
    \label{thm:hastings_site_action_distance}
    \lean{SimpleGraph.edist_smul_of_adj_iff, SimpleGraph.cycleGraph_adj_smul}
    \leanok
    If every $g\in G$ preserves the adjacency of a graph on $V$, then
    $d(g\cdot x,g\cdot y)=d(x,y)$ for all sites $x$, $y$. The rotations of a ring
    of $n$ sites are such an action.
\end{theorem}
\begin{proof}\leanok
    A walk from $x$ to $y$ is carried to a walk of the same length from $g\cdot x$ to
    $g\cdot y$, so $d(g\cdot x,g\cdot y)\le d(x,y)$; apply this to $g^{-1}$ for the
    reverse inequality.
\end{proof}
